% =====================================================================
%  1bat-ud5-4.tex  —  SESSIÓ 4: Gimnàstica de percentatges encadenats
%  (sense preàmbul: s'inclou des de main.tex)
% =====================================================================
\horatitol{Sessió 4 — Gimnàstica de percentatges encadenats}%
{Formalitzem la trampa: els canvis percentuals successius es combinen multiplicant els índexs, i això ens dóna el percentatge global equivalent.}

\subsection{Idea clau}

A la sessió 3 vau descobrir que pujar un $15\%$ i baixar un $10\%$ no torna a l'inici.
Avui formalitzem \emph{per què} i, sobretot, \emph{com} es calcula de pressa: quan
s'apliquen diversos canvis percentuals successius, els \textbf{índexs de variació es
multipliquen}.

\begin{idea}
\textbf{Percentatges encadenats}\\
Si un valor pateix diversos canvis successius amb índexs $I_1, I_2, \dots, I_n$, el
valor final és:
\[
\text{valor final}=\text{valor inicial}\per I_1\per I_2\per\cdots\per I_n.
\]
El producte $I_{\text{global}}=I_1\per I_2\per\cdots\per I_n$ és l'\textbf{índex
global}, i d'ell es llegeix el \textbf{percentatge global equivalent}.
\smallskip

\textbf{Exemple central (la llum):} pujar un $15\%$ i baixar un $10\%$ és
\[
1,15\per 0,90=1,035,
\]
és a dir, una pujada \textbf{global} del $3,5\%$ (no del $5\%$!). L'ordre dels factors
no altera el producte: el resultat és el mateix tant si puja primer com si baixa
primer.
\end{idea}

\subsection{Exemples}

\begin{exemple}[de la trampa a la fórmula]
La factura de la llum de $\num{80}\eur$ amb una pujada del $15\%$ i una baixada del
$10\%$. En comptes de fer dos passos, multipliquem els índexs d'un cop:
\[
\num{80}\per\underbrace{(1,15\per 0,90)}_{=\,1,035}=\num{80}\per 1,035=\num{82,8}\eur.
\]
Mateix resultat que a la sessió 3, però d'una sola tacada. I a més sabem el
\textbf{percentatge global}: $1,035\rightarrow +3,5\%$.
\end{exemple}

\begin{exemple}[descomptes successius i el global equivalent]
Una botiga fa un $30\%$ i després un $20\%$ addicional. L'índex global és
\[
0,70\per 0,80=0,56.
\]
Com $0,56=1-0,44$, els dos descomptes equivalen a un \textbf{únic descompte del
$44\%$} (no del $50\%$). Sobre un article de $\num{50}\eur$:
$\num{50}\per 0,56=\num{28}\eur$, exactament el que vam obtenir pas a pas.
\end{exemple}

\subsection{Exercicis resolts}

\begin{exresolt}[índex global i percentatge equivalent]
Un servei puja un $10\%$ un any i un $5\%$ l'any següent. Calcula l'índex global i
digues a quin percentatge global de pujada equival.
\solucio
Multipliquem els índexs:
\[
I_{\text{global}}=1,10\per 1,05=1,155.
\]
Com $1,155=1+0,155$, els dos augments equivalen a una pujada global del
$\mathbf{15,5\%}$ (no del $15\%$ que donaria sumar-los). Sobre, per exemple,
$\num{200}\eur$: $\num{200}\per 1,155=\num{231}\eur$.
\resposta{$I_{\text{global}}=1,155$; pujada global del $15,5\%$.}
\end{exresolt}

\begin{exresolt}[IVA sobre un preu ja rebaixat]
Un material té un preu base de $\num{400}\eur$. Primer s'hi aplica un descompte del
$25\%$ i després l'IVA del $21\%$. Calcula el preu final amb l'índex global i digues
quina variació global representa.
\solucio
Índex global:
\[
I_{\text{global}}=0,75\per 1,21=0,9075.
\]
Preu final: $\num{400}\per 0,9075=\num{363}\eur$. Com $0,9075=1-0,0925$, el resultat
global és una \textbf{baixada del $9,25\%$} respecte del preu base: el descompte del
$25\%$ pesa més que l'IVA del $21\%$.
\resposta{$\num{363}\eur$; índex global $0,9075$, equivalent a una baixada del
$9,25\%$.}
\end{exresolt}

\subsection{Exercicis per fer}

\begin{exercicis}
\begin{exlist}
  \item Calcula l'\textbf{índex global} de cada encadenament:
    \begin{enumerate}[label=\alph*), leftmargin=1.6em, topsep=2pt]
      \item pujar un $20\%$ i després un $10\%$;
      \item baixar un $30\%$ i després un $10\%$;
      \item pujar un $25\%$ i després baixar un $20\%$.
    \end{enumerate}
  \item Per a cada índex global de l'exercici 1, digues a quin \textbf{percentatge
        global} equival (pujada o baixada).
  \item Una factura de $\num{150}\eur$ puja un $12\%$ i després baixa un $5\%$. Calcula
        el preu final amb l'índex global i digues quina variació global representa.
  \item Dos descomptes seguits del $40\%$ i del $10\%$. \textbf{(Compte)} A quin
        descompte global \emph{únic} equivalen? (No és el $50\%$.)
  \item \textbf{(Interpretació)} Una nòmina puja un $2\%$ cada any durant tres anys
        seguits. Escriu l'índex global com un producte de tres factors i calcula a quin
        percentatge global equival al cap dels tres anys. Per què és més que un $6\%$?
  \item \textbf{(Raonament)} Per què, per combinar percentatges successius,
        \textbf{multipliquem} els índexs en comptes de sumar els percentatges? Relaciona
        la resposta amb la idea, de la sessió 3, que cada canvi s'aplica sobre una base
        diferent.
\end{exlist}
\end{exercicis}

% ---------------------------------------------------------------------
\fullsolucions{Sessió 4}
\begin{solucions}
\begin{sollist}
  \item (a) $1,20\per 1,10=1,32$; \quad (b) $0,70\per 0,90=0,63$; \quad
        (c) $1,25\per 0,80=1,00$.
  \item (a) $1,32=1+0,32\rightarrow$ pujada global del $32\%$; \quad
        (b) $0,63=1-0,37\rightarrow$ baixada global del $37\%$; \quad
        (c) $1,00\rightarrow$ \textbf{cap variació} (es compensen exactament; aquest
        cas és l'excepció en què sí es torna a l'inici, perquè $1,25\per 0,80=1$).
  \item $I_{\text{global}}=1,12\per 0,95=1,064$. Preu final:
        $\num{150}\per 1,064=\num{159,6}\eur$. Equival a una pujada global del $6,4\%$.
  \item $0,60\per 0,90=0,54=1-0,46$: equivalen a un \textbf{descompte únic del $46\%$}
        (no del $50\%$).
  \item $I_{\text{global}}=1,02\per 1,02\per 1,02=1,02^3\approx 1,0612$. Equival a una
        pujada global d'aproximadament el $\mathbf{6,12\%}$, una mica \textbf{més} que
        el $6\%$ de sumar $2+2+2$, perquè cada any el $2\%$ s'aplica sobre una nòmina ja
        una mica més alta (l'interès «sobre l'interès», idea que reapareixerà amb
        l'interès compost).
  \item Resposta oberta. Idea: cada canvi percentual s'aplica sobre el valor que hi ha
        \emph{en aquell moment}, que ja inclou els canvis anteriors; per això no es
        poden sumar els percentatges (que suposaria aplicar-los tots sobre la base
        original). Multiplicar els índexs respecta que cada factor actua sobre el
        resultat acumulat fins llavors.
\end{sollist}
\end{solucions}
